% ==============================================================================
% Section 2: Overview of FDI in Switzerland
% ==============================================================================

\section{Overview of FDI in Switzerland}

The aggregate foreign direct investment profile of Switzerland is characterized by extraordinary scale relative to domestic economic output. A thorough empirical assessment requires distinguishing between \textbf{capital stocks}---which measure the cumulative book or market value of direct investment assets at a specific point in time---and \textbf{capital flows} (transactions), which represent net financial movements recorded in the balance of payments over a calendar year.

\subsection{FDI Capital Stocks (Position at Year-End 2024)}

According to official Swiss National Bank statistics, Switzerland maintained the following direct investment positions at the close of 2024:

\begin{itemize}[leftmargin=1.5em, itemsep=3pt]
    \item \textbf{Inward FDI Capital Stock:} Reached \textbf{$\CHF\,925.7\bn$}, comprising:
    \begin{itemize}[leftmargin=1.2em, itemsep=1.5pt]
        \item \textit{Equity Capital:} $\CHF\,886.7\bn$ ($95.8\pct$ of total inward stock)
        \item \textit{Debt Instruments (Intra-company loans):} $\CHF\,39.0\bn$ ($4.2\pct$ of total inward stock)
    \end{itemize}
    \item \textbf{Outward FDI Capital Stock:} Reached \textbf{$\CHF\,1,340.1\bn$}, comprising:
    \begin{itemize}[leftmargin=1.2em, itemsep=1.5pt]
        \item \textit{Equity Capital:} $\CHF\,1,303.0\bn$ ($97.2\pct$ of total outward stock)
        \item \textit{Debt Instruments (Intra-company loans):} $\CHF\,37.2\bn$ ($2.8\pct$ of total outward stock)
    \end{itemize}
\end{itemize}

\begin{table}[H]
\centering
\small
\caption{Consolidated Foreign Direct Investment Capital Stocks and Flows in Switzerland (2024)}
\label{tab:fdi_overview_2024}
\renewcommand{\arraystretch}{1.2}
\begin{tabularx}{\textwidth}{>{\raggedright\arraybackslash}X r r r r r}
\toprule
\rowcolor{tableheadbg}
\textbf{Direction} & \textbf{Equity ($\CHF$\,bn)} & \textbf{Debt ($\CHF$\,bn)} & \textbf{Total Stock} & \textbf{Equity Share} & \textbf{2024 Flow} \\
\midrule
\textbf{Inward FDI (in Switzerland)} & 886.7 & 39.0 & 925.7 & 95.8\% & -83.4 \\
\rowcolor{tablerowbg}
\textbf{Outward FDI (Swiss Abroad)} & 1,303.0 & 37.2 & 1,340.1 & 97.2\% & -11.1 \\
\midrule
\textbf{Net Direct Investment Gap} & \textbf{+416.3} & \textbf{-1.8} & \textbf{+414.4} & --- & \textbf{+72.3} \\
\bottomrule
\end{tabularx}
\vspace{0.3em}
\footnotesize
\textit{Source: Swiss National Bank (SNB) Foreign Direct Investment Datasets (2024). Net Direct Investment Gap = Outward Stock minus Inward Stock. Positive gap denotes net capital export position.}
\normalsize
\end{table}

The empirical breakdown reveals that both inward and outward positions are overwhelmingly dominated by equity capital ($\ge 95.8\pct$). This high equity concentration indicates that foreign multinational enterprises operate in Switzerland with substantive capitalization and operational independence, while Swiss MNEs fund their overseas subsidiaries primarily through equity injections and retained corporate earnings rather than highly leveraged cross-border debt channels.

\subsection{FDI Capital Transactions (Flows) in 2024}

In stark contrast to the cumulative stock positions, direct investment flows in 2024 registered significant net disinvestments across both directions:

\begin{itemize}[leftmargin=1.5em, itemsep=3pt]
    \item \textbf{Inward Capital Transactions (2024):} \textbf{$\CHF\,-83.4\bn$} (Net capital withdrawal / disinvestment)
    \item \textbf{Outward Capital Transactions (2024):} \textbf{$\CHF\,-11.1\bn$} (Net foreign capital repatriation)
\end{itemize}

\begin{remarkbox}[Deconstructing the 2024 Capital Outflow Phenomenon]
The pronounced negative inward flow ($\CHF\,-83.4\bn$) in 2024 does not signify a structural collapse in Switzerland's industrial competitiveness. Rather, it reflects three extraordinary macroeconomic and institutional developments:
\begin{enumerate}[leftmargin=1.2em, itemsep=2pt]
    \item \textbf{Monetary Normalization and Interest Rate Differentials:} The conclusion of the SNB's negative interest rate policy (exit from $-0.75\pct$ policy rate) and rapid global central bank tightening prompted multinational treasuries to reallocate excess liquid reserves toward higher-yielding sovereign debt markets outside Switzerland.
    \item \textbf{Earnings Repatriation by Foreign Parent Corporations:} Heightened geopolitical volatility, energy uncertainties in Europe, and elevated US corporate borrowing costs incentivized foreign parents (particularly US-based MNEs) to repatriate accrued earnings and liquidate short-term intra-group credit balances from their Swiss finance and holding entities.
    \item \textbf{Corporate Restructuring under BEPS Pillar Two:} Global implementation of the OECD/G20 Minimum Corporate Tax rate of $15\pct$ triggered substantial organizational streamlining of Special Purpose Entities (SPEs) and non-operational holding subsidiaries in Switzerland, resulting in substantial balance sheet contractions.
\end{enumerate}
\end{remarkbox}

\begin{theorem}{Stock-Flow Accounting Dynamic in Balance of Payments}{stock_flow_accounting}
The change in direct investment capital stock between period $t-1$ and period $t$ does not equal the recorded capital transaction flow ($\text{Flow}_t$) alone. Instead, it obeys the fundamental valuation identity:
\begin{equation}
    \text{Stock}_t = \text{Stock}_{t-1} + \text{Flow}_t + \Delta \text{Valuation}_t^{\text{FX}} + \Delta \text{Valuation}_t^{\text{Price}} + \text{Other Adjustments}_t
\end{equation}
where $\Delta \text{Valuation}_t^{\text{FX}}$ denotes exchange rate revaluation effects (appreciation of the CHF decreases the domestic currency value of foreign assets), and $\Delta \text{Valuation}_t^{\text{Price}}$ captures equity price asset revaluations and write-downs.
\end{theorem}
